% Regression: p3_regions — imported current-behavior fixture from the handoff corpus.
% Formula: R = ([2,4]x[2,4]) \ {(4,4)}  — union of a 2x3 and a 3x2 injective
% p3_regions — arXiv:1804.04964 (normal PEPS FT), Sec. IV, p. 15.
% Injective-region examples before Theorem 3.
%
% Formulae:
%   R = ([2,4]x[2,4]) \ {(4,4)}  — union of a 2x3 and a 3x2 injective
%       block, hence injective (Lemma "injective_union"); the notch at
%       the removed corner site must read.
%   S = [2,4]x[2,4]              — R plus the corner site, again a union
%       of injective 2x3 / 3x2 blocks.
%   T = (whole 6x7 window) \ (staircase R' u S')  — the COMPLEMENT of a
%       staircase region: an even-odd ring (outer loop at the window
%       boundary, inner loop around the hole), injective if the PEPS is
%       at least 5x6.
%
% Divergence noted: the paper fills T opaquely, hiding the covered sites;
% tenkz's doctrine draws quiet membership tints UNDER the lattice ink, so
% member sites stay visible inside the shaded ring.
\documentclass[varwidth=24cm, border=6pt]{standalone}
\usepackage{amsmath, amssymb}
\usepackage{tenkz}
\begin{document}
% ---- R and S: a notched 3x3 block and the full 3x3 block ----
\begin{tenkzlattice}[rows=5, cols=5, boundary legs]
  \tnregion[slot=selected, label=$R$]{(2-4,2-4) - (4,4)}
\end{tenkzlattice}
\ {\rm and} \
\begin{tenkzlattice}[rows=5, cols=5, boundary legs]
  \tnregion[slot=selected, label=$S$]{(2-4,2-4)}
\end{tenkzlattice}

\bigskip
% ---- T: the ring complement of a staircase (interior hole) ----
% hole = (rows 2-4, cols 2-3) u (rows 4-5, cols 4-6)
\begin{tenkzlattice}[rows=6, cols=7, boundary legs]
  \tnregion[slot=selected, label=$T$, label at=north east]
    {(1-6,1-7) - (2-4,2-3) - (4-5,4-6)}
\end{tenkzlattice}
\end{document}
